\subsection{FURTHER PROPERTIES OF THE RADICAL}

PROPOSITION 5. Let g be a Lie algebra and r its radical.

\begin{enumerate}
\def\labelenumi{(\alph{enumi})}
\item
  If \(\rho\) is a finite-dimensional representation of \(\mathfrak{g}\) and \(\beta\) is the associated bilinear form, \(\mathfrak{r}\) and \(\mathcal{D}\mathfrak{g}\) are orthogonal with respect to \(\beta\) .
\item
  r is the orthogonal of \(\mathcal{D}\mathfrak{g}\) with respect to the Killing form.
\end{enumerate}

Let \(x, y\) be in \(\mathfrak{g}, z \in \mathfrak{r}\) . Then \([y, z] \in \mathcal{D}\mathfrak{g} \cap \mathfrak{r}\) and hence

\[
\beta ([ x, y ], z) = \beta (x, [ y, z ]) = 0
\]

(Theorem 1). Hence (a).

Let \(\mathfrak{r}'\) be the orthogonal of \(\mathcal{D}\mathfrak{g}\) with respect to the Killing form. It is an ideal of \(\mathfrak{g}\) ( \(\S 3\) , no. 6, Proposition 7 (a)) which contains \(\mathfrak{r}\) by the above. On the other hand, the image s of \(\mathfrak{r}'\) under the adjoint representation of \(\mathfrak{g}\) is solvable (Theorem 2) and hence \(\mathfrak{r}'\) is solvable being a central extension of s. Hence \(\mathfrak{r}' \subset \mathfrak{r}\) .

COROLLARY 1. Let \(\mathfrak{g}\) be a Lie algebra. Then \(\mathfrak{g}\) is solvable if and only if \(\mathcal{D}\mathfrak{g}\) is orthogonal to \(\mathfrak{g}\) with respect to the Killing form.

This is an immediate consequence of Proposition 5 (b).

COROLLARY 2. The radical \(\mathfrak{r}\) of a Lie algebra \(\mathfrak{g}\) is a characteristic ideal.

\(\mathcal{D}\mathfrak{g}\) is a characteristic ideal and the Killing form is completely invariant (§ 3, no. 6, Proposition 10). Hence the orthogonal of \(\mathcal{D}\mathfrak{g}\) with respect to the Killing form is a characteristic ideal (§ 3, no. 6, Proposition 7 (b)).

COROLLARY 3. Let g be a Lie algebra, r its radical and a an ideal of g. Then the radical of a is equal to r ∩ a.

\(r \cap a\) is a solvable ideal of \(a\) and hence is contained in the radical \(r'\) of \(a\) . Conversely, \(r'\) is an ideal of \(g\) (Corollary 2 and §1, no. 4, Proposition 2) and hence \(r' \subset r\) .

Corollary 2 can be made more precise as follows:

PROPOSITION 6. Let \(\mathfrak{g}\) be a Lie algebra, \(\mathfrak{r}\) its radical and \(\mathfrak{n}\) its largest nilpotent ideal. Every derivation of \(\mathfrak{g}\) maps \(\mathfrak{r}\) into \(\mathfrak{n}\) .

Let D be a derivation of g. Let \(g' = g + Kx_0\) be a Lie algebra in which g is an ideal of codimension 1 such that \(Dx = [x_0, x]\) for all \(x \in g\) (§ 1, no. 8, Example 1). By Corollary 3 to Proposition 5, r is contained in the radical \(r'\) of \(g'\) . Then \(D(r) = [x_0, r] \subset [g', g'] \cap r' = s'\) . For all \(x \in s' \cup g\) , \(ad_g x\) is nilpotent (Theorem 1). Hence, for all \(x \in s' \cap g\) , \(ad_g x\) is nilpotent. Hence \(D(r)\) is contained in the nilpotent ideal \(s' \cap g\) of g.

COROLLARY. The largest nilpotent ideal of a Lie algebra is a characteristic ideal.

Remark. To summarize some of the above results, note that, if r, n, s, t denote respectively the radical of g, the largest nilpotent ideal of g, the nilpotent radical of g and the orthogonal of g with respect to the Killing form, then

\[
\mathfrak {r} \supset \mathfrak {k} \supset \mathfrak {n} \supset \mathfrak {s}.
\]

The inclusion \(\mathfrak{r} \supset \mathfrak{t}\) follows from Proposition 5 (b). The inclusion \(\mathfrak{t} \supset \mathfrak{n}\) follows from §4, no. 4, Proposition 6 (b). The inclusion \(\mathfrak{n} \supset \mathfrak{s}\) has been pointed out in Remark 2 of no. 3.

